\documentclass[11pt]{article}
\usepackage[margin=1in]{geometry}
\usepackage{amsmath,amssymb}
\usepackage[hidelinks]{hyperref}
\emergencystretch=3em

\title{Literature Status: Subspaces of Banach Lattices Without Bibasic Sequences}
\author{arXiv:1907.07589 answered by arXiv:2210.00614}
\date{2026-06-29}

\begin{document}
\maketitle

\paragraph{Status.}
\textbf{literature\_already\_answered.}  This packet records a later-paper answer, not a new proof from this run.

\paragraph{Original source.}
M. A. Taylor and V. G. Troitsky, \emph{Bibasic sequences in Banach lattices}, arXiv:1907.07589.
In Remark 4.4, after showing that every closed infinite-dimensional subspace of the closed span of a bibasic sequence contains a bibasic sequence, the authors ask whether every closed infinite-dimensional subspace of a Banach lattice contains a bibasic sequence.

\paragraph{Supporting answer.}
T. Oikhberg, M. A. Taylor, P. Tradacete, and V. G. Troitsky, \emph{Free Banach lattices}, arXiv:2210.00614.
Its introduction says that Section 7 constructs the first example of a subspace of a Banach lattice without a bibasic sequence and that this answers a question from Taylor--Troitsky.  The decisive statement is Theorem 7.5:
\[
  \phi(c_0)\subseteq FBL^{(p)}[c_0]
  \quad\text{does not contain any bibasic sequence whenever }1\le p<\infty.
\]
For \(p=1\), \(FBL^{(1)}[c_0]\) is the ordinary free Banach lattice \(FBL[c_0]\).

\paragraph{Why this answers the target.}
The map \(\phi:c_0\to FBL[c_0]\) is the canonical linear isometric embedding, so \(\phi(c_0)\) is a closed infinite-dimensional subspace of the Banach lattice \(FBL[c_0]\).  Theorem 7.5 says this subspace contains no bibasic sequence.  Hence the answer to Remark 4.4 of arXiv:1907.07589 is negative.

\paragraph{Corroborating pointer.}
E. Garcia-Sanchez, D. de Hevia, and P. Tradacete, \emph{Free objects in Banach space theory}, arXiv:2302.10807, explicitly states that arXiv:1907.07589 Remark 4.4 asked whether there exists a subspace of a Banach lattice with no bibasic sequence, and that arXiv:2210.00614 Theorem 7.5 gives such a subspace by using the canonical copy of \(c_0\) in \(FBL[c_0]\).

\paragraph{Scope and limitations.}
This settles the general closed-subspace bibasic question from Remark 4.4.  It does not address the separate uo-bibasic or order-continuous variants posed later in arXiv:1907.07589.

\paragraph{Search evidence.}
The run indexes were searched for arXiv:1907.07589, bibasic sequences, Remark 4.4, subspaces without bibasic sequences, and arXiv:2210.00614.  No prior packet for this Taylor--Troitsky question was found.  The local parsed sources for arXiv:2210.00614 and arXiv:2302.10807 contain the exact cross-references quoted above.

\end{document}
